%% chapter 5 dataset, network structure, experiment and result
\chapter{讨论}
\label{cha:discussion}


\section{结果与先验预期的一致性}

从定性层面来看，本文的实验观测与既有 WebAssembly 性能研究中的主要结论总体一致。首先，在 MicroBenchmark 中，涉及高频系统接口调用与宿主边界穿越的程序普遍表现出较大的 \path{native-better} 比值，这与 WASI 场景下系统接口实现质量对端到端性能产生显著影响的已有结论相一致\cite{mao2024ewapa}。其次，在部分大规模数值计算 kernel（如 \path{gemm}）中观察到 Wasm 相对劣势显著的现象，也与已有研究中关于“小规模 kernel 上的性能结论不宜直接推广至复杂 workload”的讨论相呼应\cite{jangda2021not}。

与此同时，本文在部分小规模计算密集型微基准中观察到 \path{wasm-better} 或 \path{similar} 的结果，表明“计算密集”这一单一特征不足以作为性能优劣的充分判据。程序的热点规模、循环结构以及后端优化空间等因素同样对性能结果产生重要影响，这一点在第 4.2 节与第 4.3 节的分析中已有体现。

\section{与相关工作的关系}

相较于侧重端到端运行时间或浏览器执行环境的研究\cite{sunarto2023systematic,yan2021understanding}，本文在 Non-Web 场景下，基于 WASI 与单一运行时环境开展实验，并将性能差异与 LLVM IR 层面的静态结构特征相结合，从而为理解“何种程序结构更易放大 Wasm 与 native 的性能差距”提供了新的分析视角。

相较于面向大规模 benchmark 的对比研究\cite{jangda2021not}，本文在样本规模上有所限制，但通过构建结构可控的微基准以隔离单一因素，并采用与热点区域对齐的内部计时方法，使静态特征与性能指标在语义上具有更强的一致性。需要指出的是，本文未对多种 Wasm 运行时进行系统比较，因此相关结论在不同运行时实现及工具链配置下的适用性仍有待进一步验证。

\section{实践意义与适用边界}

从工程实践角度来看，本文结果可归纳为若干具有参考价值的经验性结论。首先，当程序核心路径较短且频繁涉及宿主接口调用，或动态内存分配占比较高，或存在大量难以预测的不规则分支结构时，Wasm 相对于 native 出现显著性能劣势的可能性较大。其次，当程序以规则数值循环为主且更接近内存带宽受限（memory-bound）特征时，Wasm 与 native 之间的性能差距可能趋于收敛，但具体表现仍需通过实际测量加以验证。

此外，本文所构建的静态特征体系及分类模型可作为开发早期的粗粒度筛查工具，用于初步评估程序在 Wasm 环境下的潜在性能表现。然而，该方法不应替代在目标输入规模、具体运行时及硬件环境下的性能测试与分析。需要强调的是，本文未涉及 WebAssembly 在安全隔离、可移植性及部署灵活性等方面的非性能优势；在相关需求权重较高且性能差距可接受的情况下，Wasm 仍可能构成合理的工程选择。

\section{研究局限}

除第 4.4 节所述的样本规模与标签构造问题外，本文还存在以下若干限制。首先，静态特征依赖于特定编译阶段与 LLVM IR 表示，不同优化等级或前端编译策略可能改变特征分布及其与性能之间的关系。其次，实验在单一容器化 Linux 环境中完成，未系统考察操作系统版本、CPU 微架构及功耗管理策略等因素的影响。最后，性能测量基于内部计时方法，尚未结合硬件性能计数器对执行瓶颈进行细粒度分解。此外，尽管通过 warmup 缓解了 JIT 启动开销的影响，但实验结果仍主要反映短期稳态行为，与长期运行场景可能存在差异。

因此，本文结论应理解为在受控实验条件下得到的结构性观察，而非对所有 WebAssembly 应用场景的普适性结论。



